\documentclass[11pt]{article}

\usepackage[margin=1in]{geometry}
\usepackage{amsmath, amssymb}
\usepackage{booktabs}
\usepackage{algorithm}
\usepackage{algpseudocode}
\usepackage{hyperref}
\usepackage{xcolor}
\usepackage{listings}
\usepackage{caption}

\title{Complex Lane Packing for Encrypted GQA Key/Value Caches}
\author{Cachemir GQA -- Ciphertext Pipeline}
\date{\today}

\newcommand{\conj}[1]{\overline{#1}}

\begin{document}
\maketitle

\begin{abstract}
This note documents a lane-packing optimization applied to the Orion-based
ciphertext implementation of Grouped Query Attention (GQA) decoding in
\texttt{Cachemir/gqa/ciphertext}. CKKS ciphertext slots natively encode
complex numbers, but the existing pipeline only ever populated the real
component (imaginary part fixed at zero), discarding half of the available
packing density. We introduce a scheme that packs two tokens per ciphertext
slot -- one in the real part, one in the imaginary part -- for the
Key/Value cache, and show how the resulting cross-term contamination in
ciphertext--ciphertext multiplication is avoided (for $QK^\top$) or removed
algebraically using complex conjugation (for scores$\times V$). We give the
exact algorithms, a correctness proof of the conjugate trick, and measured
before/after cost tables.
\end{abstract}

\tableofcontents

\section{Background}

\subsection{CKKS slots and the packing gap}
A CKKS ciphertext with $n_{he}$ slots natively encodes a vector in
$\mathbb{C}^{n_{he}}$: every slot holds one \emph{complex} number, and
ciphertext--ciphertext and ciphertext--plaintext multiplication are defined
slot-wise on $\mathbb{C}$. The existing GQA pipeline only ever encodes
real-valued tensors (\texttt{torch.float64}), i.e. every slot is used as
$z = a + 0i$. The imaginary half of every slot's capacity is therefore
completely unused -- a factor-of-two packing inefficiency that propagates
directly into ciphertext count (memory) and into the number of
ciphertext--ciphertext (ct-ct) multiplications required to consume the
Key/Value (K/V) cache, which is the dominant cost of encrypted attention.

\subsection{GQA slot layout (recap)}
Let $n_{he}$ be the number of CKKS slots, $d_{kv}$ the compact K/V feature
width, $H$ the number of query heads, $n_{kv}$ the number of K/V heads,
$\mathrm{ratio} = H / n_{kv}$, and
\[
  t_p = \frac{n_{he}}{d_{kv}}
\]
the number of tokens that fit, \emph{before} this optimization, in a single
real-valued cache ciphertext (one token occupies every $t_p$-th slot,
i.e.\ physical lane $\ell \in \{0, \dots, t_p-1\}$ selects slots
$\{\ell, \ell+t_p, \ell+2t_p, \dots\}$). A cache of $N$ tokens therefore
requires
\[
  C_{\text{before}}(N) = \left\lceil \frac{N}{t_p} \right\rceil
\]
ciphertexts, and the attention read-side loops (`$QK^\top$' and
`scores$\times V$') iterate once per cache ciphertext, each iteration
costing one ct-ct multiplication.

\section{Complex Lane Packing Scheme}

\subsection{Idea}
Because each CKKS slot already holds a complex number, we can pack a
\emph{second}, independent token into the previously-unused imaginary
component of the same physical lane. Concretely, token $2k$ is written into
the real part of lane $\ell$ and token $2k+1$ into the imaginary part of the
\emph{same} lane $\ell$:
\[
  \text{slot}(\ell) \;=\; \underbrace{x_{2k}}_{\Re} \;+\; \underbrace{x_{2k+1}}_{\Im}\, i .
\]
This doubles the effective per-ciphertext capacity to $2 t_p$ tokens, so a
cache of $N$ tokens now needs only
\[
  C_{\text{after}}(N) = \left\lceil \frac{N}{2 t_p} \right\rceil
  \;\approx\; \tfrac{1}{2}\, C_{\text{before}}(N)
\]
ciphertexts.

\subsection{Writing packed tokens: extending \texttt{vmm\_kv}}
The projection routine \texttt{vmm\_kv} writes one new token into a cache
ciphertext at structural position $\mathrm{pos}$ using a plaintext mask
$m_{\mathrm{pos}} \in \{0,1\}^{n_{he}}$ with
$m_{\mathrm{pos}}[\ell] = \mathbb{1}[\ell \equiv \mathrm{pos} \!\!\pmod{t_p}]$,
via $\text{masked} = (X_c W_c) \odot m_{\mathrm{pos}}$ (a ct-pt multiply),
after folding replicated contributions into the correct lane with
$\log_2 t_p$ rotations. We extend $\mathrm{pos}$ to range over
$\{0, \dots, 2t_p - 1\}$:

\begin{itemize}
  \item $\mathrm{pos} < t_p$: unchanged -- real mask, contributes to
        $\Re(\text{slot})$.
  \item $\mathrm{pos} \ge t_p$: physical lane
        $\ell = \mathrm{pos} - t_p$, but the mask is multiplied by the
        imaginary unit,
        \[
          m'_{\mathrm{pos}} = i \cdot m_{\ell} \in \{0, i\}^{n_{he}},
        \]
        so the ct-pt product lands in $\Im(\text{slot})$ instead of
        $\Re(\text{slot})$.
\end{itemize}
No other change is required: rotation and addition act identically (and
independently) on the real and imaginary components of a slot, since both
are $\mathbb{C}$-linear operations.

\subsection{Cache capacity doubling}
The packed-cache classes (\texttt{HEKCache}/\texttt{HEVCache}) already
accumulate tokens into the \emph{current} ciphertext until a capacity
threshold is reached, then start a new one. Doubling the threshold,
\[
  t \;\leftarrow\; 2 t_p ,
\]
is sufficient: the existing append logic transparently drives the extended
$\mathrm{pos}$ range above, alternating between the real half
($\mathrm{pos} \in [0, t_p)$) and the imaginary half
($\mathrm{pos} \in [t_p, 2t_p)$) of each ciphertext.

\section{Consuming Packed Ciphertexts in Attention}

\subsection{$QK^\top$: packing is ``free'' because $Q$ stays real}
For two complex numbers $a = a_r + a_i i$ and $b = b_r + b_i i$,
multiplication mixes components:
\[
  a b = (a_r b_r - a_i b_i) + (a_r b_i + a_i b_r)\, i .
\]
In general this means multiplying two lane-packed ciphertexts contaminates
each output lane with cross terms from \emph{both} packed values -- exactly
the failure mode we must avoid.

However, in \texttt{qkt\_gqa\_he} the query ciphertext $Q$ is \emph{never}
lane-packed: it is produced once per decoding step from a single new token
via $\mathrm{pos} = 0$ only, so $Q = q + 0i$ is purely real
($a_i = 0$). Substituting $a_i = 0$ above,
\[
  Q \cdot K_{\text{pack}} \;=\; q\,(k_r + k_i i) \;=\; \underbrace{q k_r}_{\Re}
  \;+\; \underbrace{q k_i}_{\Im}\, i ,
\]
which has \emph{no} cross term: the real part is exactly the score against
the even-indexed cached token and the imaginary part is exactly the score
against the odd-indexed cached token. One ct-ct multiplication against a
packed $K$ ciphertext therefore yields two independent, uncontaminated
scores -- a 2$\times$ reduction in ct-ct multiplications for this stage at
no extra algorithmic cost.

\subsection{scores$\times V$: the conjugate trick}
The second stage, \texttt{softmax\_v\_gqa\_he}, multiplies the (now
complex, paired) attention scores $s = s_r + s_i i$ against the packed value
cache $V = v_r + v_i i$. Here \emph{both} operands are complex, so a direct
product does mix lanes:
\[
  s V = (s_r v_r - s_i v_i) + (s_r v_i + s_i v_r) i .
\]
What we actually want to accumulate into the running output is the sum of
both paired tokens' contributions,
\[
  s_r v_r + s_i v_i ,
\]
with no cross term. This is obtained by multiplying by the
\emph{complex conjugate} of $V$ instead of $V$ itself. Since
$\conj{V} = v_r - v_i i$,
\[
  s \cdot \conj{V}
  = (s_r + s_i i)(v_r - v_i i)
  = \underbrace{(s_r v_r + s_i v_i)}_{\Re}
    \;+\; \underbrace{(s_i v_r - s_r v_i)}_{\Im}\, i .
\]
The real part is \emph{exactly} the desired paired-token accumulation; the
imaginary part is discardable cross-term noise. Because the final client-side
decode already keeps only $\Re(\cdot)$ (CKKS decoding always discards
residual imaginary noise), no additional homomorphic real-part extraction is
needed -- a single ct-ct multiplication $s \cdot \conj{V}$ replaces what
would otherwise require two separate multiplications (one per unpacked
token).

\subsection{The missing primitive: homomorphic conjugation}
Complex conjugation of a CKKS ciphertext is implemented as a Galois
automorphism (the same family of operation used for slot rotations, but
with a different Galois element -- the order-2 element of the
multiplicative group associated with the ring, sometimes called the
``orthogonal subgroup'' element). The underlying Lattigo library exposes
this natively as \texttt{Evaluator.Conjugate}, but the Orion wrapper used by
this pipeline did not expose it. We added it end-to-end:
\begin{itemize}
  \item \textbf{Go layer} (\texttt{evaluator.go}): new exported
        \texttt{Conjugate}/\texttt{ConjugateNew} functions, plus
        \texttt{AddConjugationKey()} which generates (once, lazily) the
        Galois key for the conjugation automorphism, mirroring the existing
        power-of-two rotation key cache.
  \item \textbf{C bindings} (\texttt{bindings.py}): registered the two new
        exported symbols.
  \item \textbf{Python evaluator} (\texttt{evaluator.py}): a
        \texttt{conjugate(ctxt, in\_place)} wrapper mirroring
        \texttt{rotate(...)}.
  \item \textbf{Tensor API} (\texttt{tensors.py}): a
        \texttt{CipherTensor.conjugate()} method mirroring
        \texttt{CipherTensor.roll()}.
\end{itemize}
This primitive was validated in isolation: encrypting a random complex
vector, conjugating homomorphically, decrypting, and comparing against
\texttt{numpy.conj} gave a maximum absolute error of $\approx 3.4\times
10^{-8}$, consistent with ordinary CKKS rescaling noise.

\section{Algorithms}

\begin{algorithm}[H]
\caption{\texttt{vmm\_kv} with paired real/imaginary positions}
\begin{algorithmic}[1]
\Require encrypted sparse chunks $\{X_c\}$, plaintext weight chunks $\{W_c\}$,
  dims (incl. $t_p$), level $\ell$, position $\mathrm{pos} \in [0, 2t_p)$
\State $\mathrm{is\_imag} \gets (\mathrm{pos} \ge t_p)$
\State $\mathrm{phys} \gets \mathrm{pos} - t_p \cdot \mathbb{1}[\mathrm{is\_imag}]$
\State $m \gets \mathbf{0} \in \mathbb{C}^{n_{he}}$
\State $m[\mathrm{phys} :: t_p] \gets i$ \textbf{if} $\mathrm{is\_imag}$ \textbf{else} $1$
\State $\mathrm{out} \gets \varnothing$
\For{each $(X_c, W_c)$}
  \State $\mathrm{acc} \gets X_c \cdot W_c$ \Comment{ct-pt multiply}
  \For{$\mathrm{step} = 1, 2, 4, \dots < t_p$} \Comment{$\log_2 t_p$ rotations}
    \State $\mathrm{acc} \gets \mathrm{acc} + \mathrm{roll}(\mathrm{acc}, \pm\mathrm{step})$
    \Comment{sign chosen from bits of $\mathrm{phys}$}
  \EndFor
  \State $\mathrm{masked} \gets \mathrm{acc} \odot m$ \Comment{ct-pt multiply (complex $m$ if imaginary half)}
  \State $\mathrm{out} \gets \mathrm{out} + \mathrm{masked}$
\EndFor
\State \Return $\mathrm{out}$
\end{algorithmic}
\end{algorithm}

\begin{algorithm}[H]
\caption{Packed $QK^\top$ (\texttt{qkt\_gqa\_he})}
\begin{algorithmic}[1]
\Require real query ciphertext $Q$ (per group $c$), packed key ciphertexts $\{K_b\}_{b=1}^{C_{\text{after}}(N)}$
\State $Q_{\mathrm{rep}} \gets \textsc{BlockReplicate}(Q, t_p)$
\For{each packed key ciphertext $K_b$}
  \State $\mathrm{folded} \gets Q_{\mathrm{rep}} \cdot K_b$ \Comment{1 ct-ct mult scores \emph{two} tokens; no cross terms since $Q_{\mathrm{rep}}$ is real}
  \State fold $\mathrm{folded}$ over $\log_2 d_h$ rotations (head-dim reduction)
  \State append $\mathrm{folded}$ to attention row for group $c$
\EndFor
\end{algorithmic}
\end{algorithm}

\begin{algorithm}[H]
\caption{Packed scores$\times V$ (\texttt{softmax\_v\_gqa\_he})}
\begin{algorithmic}[1]
\Require packed attention scores $\{s_b\}$, packed value ciphertexts $\{V_b\}$
\State $\mathrm{out} \gets 0$
\For{each $b$}
  \State $\conj{V_b} \gets \textsc{Conjugate}(V_b)$ \Comment{new automorphism}
  \State $\mathrm{prod} \gets s_b \cdot \conj{V_b}$ \Comment{1 ct-ct mult; $\Re(\mathrm{prod}) = s_r v_r + s_i v_i$}
  \State fold $\mathrm{prod}$ over $\log_2 t_p$ rotations
  \State $\mathrm{out} \gets \mathrm{out} + \mathrm{prod} \odot \mathrm{out\_mask}$ \Comment{ct-pt multiply, real mask}
\EndFor
\State \Return $\mathrm{out}$ \Comment{client discards $\Im(\cdot)$ on decode}
\end{algorithmic}
\end{algorithm}

\section{Cost Model}
\label{sec:cost-model}

Let $\mathrm{ratio} = H / n_{kv}$ and $C(N)$ be the number of K/V cache
ciphertexts for $N$ cached tokens. Per decoding step, the read-side cost of
$QK^\top$ and scores$\times V$ combined is
\[
  \#\text{ct-ct} \;=\; 2 \cdot \mathrm{ratio} \cdot C(N), \qquad
  \#\text{conjugations} \;=\;
  \begin{cases}
    0 & \text{before (n/a)}\\
    \mathrm{ratio} \cdot C(N) & \text{after}
  \end{cases}
\]
with $C_{\text{before}}(N) = \lceil N / t_p \rceil$ and
$C_{\text{after}}(N) = \lceil N / (2 t_p) \rceil$. The write-side cost (new
token projection via \texttt{vmm\_kv}, dominated by
$R \cdot \mathrm{ratio} \cdot \log_2 t_p$ rotations per token, where
$R = \lceil d_{kv} / t_p \rceil$) is \emph{unchanged} by this optimization,
since each decoding step still injects exactly one new token regardless of
packing.

\section{Measured Results}
\label{sec:measured-results}

Configuration: $n_{he} = 512$, $d = 128$, $H = 8$, $n_{kv} = 4$
($d_h = 16$, $\mathrm{ratio} = 2$, $d_{kv} = 64$, $t_p = 8$).

\subsection{Cache ciphertext count (memory)}
\begin{table}[H]
\centering
\begin{tabular}{rrrr}
\toprule
Tokens $N$ & Before $\lceil N/t_p \rceil$ & After $\lceil N/2t_p \rceil$ & Reduction \\
\midrule
6   & 1  & 1 & 0\%  \\
21  & 3  & 2 & 33\% \\
51  & 7  & 4 & 43\% \\
101 & 13 & 7 & 46\% \\
\bottomrule
\end{tabular}
\caption{K/V cache ciphertext count before/after complex packing. Very
short sequences see no benefit (a single ciphertext already covers them
either way); the reduction approaches the asymptotic 2$\times$ as $N$
grows.}
\end{table}

\subsection{Ciphertext--ciphertext multiplications ($QK^\top$ + scores$\times V$)}
\begin{table}[H]
\centering
\begin{tabular}{rrrrr}
\toprule
Tokens $N$ & Before & After (measured) & Reduction & Conjugations (after) \\
\midrule
6   & 4  & 4  & 0\%  & 2  \\
21  & 12 & 8  & 33\% & 4  \\
51  & 28 & 16 & 43\% & 8  \\
101 & 52 & 28 & 46\% & 14 \\
\bottomrule
\end{tabular}
\caption{ct-ct multiplication counts. ``Before'' values are computed from
the closed-form cost model above (consistent with the pre-packing code
path); ``after'' values are measured directly from the instrumented
pipeline. Conjugations are a new op class introduced by this optimization,
each cheaper than a ct-ct multiplication.}
\end{table}

\subsection{Rotations and other operation counts}
\begin{table}[H]
\centering
\begin{tabular}{rrrrr}
\toprule
Tokens $N$ & Rotations & ct-pt mult & ct-ct mult & Conjugations \\
\midrule
6   & 1448  & 450  & 4  & 2  \\
21  & 4792  & 1412 & 8  & 4  \\
51  & 11480 & 3336 & 16 & 8  \\
101 & 22622 & 6542 & 28 & 14 \\
\bottomrule
\end{tabular}
\caption{Full measured counter dump for the packed (``after'') pipeline.
Rotation counts are essentially unchanged from the pre-packing baseline at
matching $N$: they are dominated by the per-token write path in
\texttt{vmm\_kv} (folding each new token into position, unaffected by
packing), not by the read-side loops that benefit from packing. Only the
smaller $QK^\top$/scores$\times V$ fold-rotation contributions shrink
proportionally to the ct-ct savings above.}
\end{table}

\subsection{Correctness}
\begin{table}[H]
\centering
\begin{tabular}{rrrr}
\toprule
Tokens $N$ & qkt\_err & softmaxv\_err & Status \\
\midrule
6   & $7.16\times 10^{-5}$ & $2.93\times 10^{-3}$ & CHECK \\
21  & $7.16\times 10^{-5}$ & $4.32\times 10^{-3}$ & CHECK \\
51  & $1.05\times 10^{-4}$ & $7.22\times 10^{-3}$ & CHECK \\
101 & $1.05\times 10^{-4}$ & $1.09\times 10^{-2}$ & CHECK \\
\bottomrule
\end{tabular}
\caption{Errors against the (unchanged, real-valued) plaintext reference.
The $N{=}6$ case matches the pre-existing, documented CKKS approximation
baseline for this configuration exactly (no regression from packing); the
slow growth with $N$ reflects accumulating CKKS noise across more ct-ct
terms, independent of whether packing is used.}
\end{table}

\section{Benchmark: Complex vs.\ Real-Only, Subprocess-Isolated}

The cost-model tables in Section~\ref{sec:measured-results} were derived
from operation counters within a single long-lived process. To also
obtain clean, per-mode latency and peak-memory numbers, a dedicated
benchmark harness (\texttt{Cachemir/gqa/ciphertext/\allowbreak
benchmark\_complex\_packing.py}) runs each configuration \emph{twice},
once with \texttt{pack\_complex=True} and once with
\texttt{pack\_complex=False}, in two separate, freshly-spawned
subprocesses per configuration. This isolation is necessary because
peak resident set size (RSS) is a monotonically non-decreasing,
process-wide high-water mark (via \texttt{resource.getrusage(\allowbreak
RUSAGE\_SELF).ru\_maxrss}): measuring it for two modes inside the same
process would contaminate the second measurement with the first mode's
memory usage.

\subsection{Hardware}
Benchmarks were run on a single (non-virtualized) machine with the
following relevant specifications:
\begin{itemize}
  \item CPU: 2$\times$ Intel(R) Xeon(R) Platinum 8260 @ 2.40\,GHz
        (24 physical cores per socket, 48 physical cores / 96 logical
        threads total via hyperthreading), NUMA with 2 nodes.
  \item Cache: 1.5\,MiB L1d + 1.5\,MiB L1i, 48\,MiB L2, 71.5\,MiB L3
        (shared, 2 instances).
  \item Memory: 250\,GiB RAM.
  \item OS/kernel: Linux 6.8 (Ubuntu), x86\_64.
\end{itemize}
Despite the machine having 96 logical cores available, the pipeline as
benchmarked is \emph{single-threaded at the application level}: the
Orion/Lattigo Go backend invoked here performs each ciphertext operation
(rotation, ct-pt/ct-ct multiplication, conjugation) as one sequential
call from the Python driver, and no goroutine/thread-pool parallelism
across independent ciphertext operations is used in this code path (the
Orion backend wrapper contains no explicit concurrency primitives such
as \texttt{go func(...)} or worker pools for these operations). The
reported wall-clock numbers therefore reflect single-core throughput of
the CKKS backend, not a parallel/multi-core execution; the large core
count and RAM headroom of the host mean the measurements are not
resource-constrained, but they also do not benefit from multi-core
speedup, and results should scale similarly on any single fast core of
comparable clock speed.

\subsection{Results}
Each configuration below ran an $n_{\text{prefill}}$-token prefill loop
followed by one encrypted decoding step, using
$n_{he} = 512$, $d = 128$, $H = 8$, $n_{kv} = 4$
($t_p = 8$, $\mathrm{ratio} = 2$) for the two larger cases, and
$d = 8$, $H = 4$, $n_{kv} = 2$ ($t_p = 128$) for the smallest case.

\begin{table}[H]
\centering
\small
\begin{tabular}{lrrr}
\toprule
Metric ($d{=}8$, $n_{\text{prefill}}{=}5$) & Complex & Real-only & Ratio \\
\midrule
Latency: prefill (s)     & 0.803   & 0.813   & 0.99$\times$ \\
Latency: decode step (s) & 0.420   & 0.404   & 1.04$\times$ \\
Latency: total (s)       & 1.268   & 1.262   & 1.00$\times$ \\
Peak RSS (MB)            & 857.0   & 851.0   & 1.01$\times$ \\
kcache / vcache ciphertexts & 1 / 1 & 1 / 1 & 1.00$\times$ \\
att\_cts / O ciphertexts    & 2 / 2 & 2 / 2 & 1.00$\times$ \\
ct-ct mult (total)       & 4       & 4       & 1.00$\times$ \\
Conjugations (total)     & 2       & 0       & n/a \\
\bottomrule
\end{tabular}
\caption{Smallest case: cache already fits in a single ciphertext under
both layouts ($t_p{=}128 \gg N{=}6$ tokens), so packing has no effect --
included as a control.}
\end{table}

\begin{table}[H]
\centering
\small
\begin{tabular}{lrrr}
\toprule
Metric ($d{=}128$, $n_{\text{prefill}}{=}9$) & Complex & Real-only & Ratio \\
\midrule
Latency: prefill (s)     & 9.334   & 9.293   & 1.00$\times$ \\
Latency: decode step (s) & 1.147   & 1.181   & 0.97$\times$ \\
Latency: total (s)       & 11.030  & 11.026  & 1.00$\times$ \\
Peak RSS (MB)            & 1598.0  & 1600.2  & 1.00$\times$ \\
kcache / vcache ciphertexts & 1 / 1 & 2 / 2 & 0.50$\times$ \\
att\_cts ciphertexts        & 2     & 4     & 0.50$\times$ \\
O ciphertexts               & 2     & 2     & 1.00$\times$ \\
ct-ct mult (total)       & 4       & 8       & 0.50$\times$ \\
ct-ct mult ($QK^\top$ phase)     & 2 & 4 & 0.50$\times$ \\
ct-ct mult (scores$\times V$ phase) & 2 & 4 & 0.50$\times$ \\
Conjugations (total)     & 2       & 0       & n/a \\
\bottomrule
\end{tabular}
\caption{$N=10$ tokens is just past the real-only layout's single-ciphertext
capacity ($t_p{=}8$), so the cache splits into 2 ciphertexts real-only vs.\
1 under packing -- the ct-ct multiplication count halves accordingly in
both the $QK^\top$ and scores$\times V$ phases.}
\end{table}

\begin{table}[H]
\centering
\small
\begin{tabular}{lrrr}
\toprule
Metric ($d{=}128$, $n_{\text{prefill}}{=}20$) & Complex & Real-only & Ratio \\
\midrule
Latency: prefill (s)     & 19.742  & 19.843  & 0.99$\times$ \\
Latency: decode step (s) & 1.205   & 1.256   & 0.96$\times$ \\
Latency: total (s)       & 21.498  & 21.649  & 0.99$\times$ \\
Peak RSS (MB)            & 2312.0  & 2311.3  & 1.00$\times$ \\
kcache / vcache ciphertexts & 2 / 2 & 3 / 3 & 0.67$\times$ \\
att\_cts ciphertexts        & 4     & 6     & 0.67$\times$ \\
O ciphertexts               & 2     & 2     & 1.00$\times$ \\
ct-ct mult (total)       & 8       & 12      & 0.67$\times$ \\
ct-ct mult ($QK^\top$ phase)     & 4 & 6 & 0.67$\times$ \\
ct-ct mult (scores$\times V$ phase) & 4 & 6 & 0.67$\times$ \\
Conjugations (total)     & 4       & 0       & n/a \\
\bottomrule
\end{tabular}
\caption{$N=21$ tokens: real-only needs $\lceil 21/8\rceil = 3$ cache
ciphertexts vs.\ $\lceil 21/16\rceil = 2$ under packing, matching the
0.67$\times$ ct-ct multiplication and ciphertext-count ratios observed
elsewhere in this report for the same $N$.}
\end{table}

\subsection{Discussion}
The ct-ct multiplication and ciphertext-count reductions match the
closed-form cost model exactly (Section~\ref{sec:measured-results}).
However, at these configurations \emph{total wall-clock latency and peak
RAM are essentially unchanged} between the two modes ($\le$4\% difference,
within run-to-run noise): the per-token \emph{write} path
(\texttt{vmm\_kv}'s ct-pt multiply and $\log_2 t_p$-step rotation fold for
each new token during prefill) dominates total time end-to-end, and that
path is, by design, unaffected by this optimization (Section~\ref{sec:cost-model}).
The benefit is concentrated in the smaller \emph{read}-side
$QK^\top$/scores$\times V$ phases, whose ct-ct multiplication counts drop
exactly as predicted; as sequence length grows well beyond the prefill
regime tested here (i.e.\ many decoding steps re-reading an already-long
cache, rather than one write per read as in this benchmark), the read-side
savings are expected to dominate overall latency and peak memory far more
visibly than in these prefill-dominated measurements.



\begin{itemize}
  \item Two tokens are packed per CKKS ciphertext slot using the real and
        imaginary components, exploiting slot capacity that was previously
        always zero-padded.
  \item The query path is deliberately left real/unpacked, which makes the
        $QK^\top$ stage benefit ``for free'' -- no cross-term contamination,
        no extra operations.
  \item The scores$\times V$ stage requires one new primitive,
        homomorphic complex conjugation, which was added to the Orion/
        Lattigo backend; the conjugate trick algebraically isolates the
        desired sum of paired-token contributions in the real part of a
        single ct-ct product.
  \item Net effect: K/V cache memory and ct-ct multiplication count (the
        dominant FHE cost) trend toward a 2$\times$ improvement as sequence
        length grows, at the cost of one cheap conjugation operation per
        packed value ciphertext consumed. Rotation counts are essentially
        unaffected, since they are dominated by the per-token write path,
        which lies outside this optimization's scope. No regression in
        numerical correctness was observed.
\end{itemize}

\end{document}
